\section{Use the pooled information when choosing the next audit}
\label{exec17:s7-constraint-design}

The earlier allocation deliberately estimated both stratum moments without
using the exact pooled audit. That is a useful baseline, but the bank should
also ask whether information it already owns changes the next purchase.
Retain precisely the same hypothetical target, the same independent truthful
neighborhood law and the same known response and monetary inputs. Grant the
pooled constraint in \eqref{exec17:s7-pooled-constraint}; no additional outcome
experiment or homogeneous graph restriction is introduced. In this section
the pooled information is already available to every compared method. Its
original acquisition cost is outside the incremental comparison, and would
have to be restored in a complete investment case.

Let $\widehat\theta_z$ be independent stratum audit means. Add any multiple
$\lambda$ of the known zero constraint to the direct value estimator:
\begin{equation}
 \begin{split}
  \widehat v_\lambda
   &=22-\tfrac\lambda4+a_0\widehat\theta_0+a_1\widehat\theta_1,\\
  a_0&=-8+\tfrac34\lambda,\qquad
  a_1=-48+\tfrac14\lambda,\qquad a_0-3a_1=136.
 \end{split}
 \label{exec17:s7-constrained-estimator}
\end{equation}
The estimator is unbiased for every feasible pair of graph moments because
$3\theta_0/4+\theta_1/4=1/4$. This is not an arbitrary change in the target:
all coefficient choices have the same expectation on the supplied information
set. The particular choice $\lambda=0$ recovers the earlier two-mean estimator.
The estimator need not itself lie in the sharp value set in a finite sample;
unbiasedness and feasibility of a realized estimate are distinct properties.

For fixed $\lambda$ and positive audit sizes its variance obeys
\begin{equation}
 \Var(\widehat v_\lambda)
 =\frac{a_0^2\theta_0(1-\theta_0)}{n_0}
  +\frac{a_1^2\theta_1(1-\theta_1)}{n_1}
 \leq \frac14\left(\frac{a_0^2}{n_0}+\frac{a_1^2}{n_1}\right).
 \label{exec17:s7-constrained-variance}
\end{equation}
Use this common one-quarter Bernoulli envelope as a declared design criterion.
It is conservative: the pooled restriction also limits which two actual
variances can occur together. Minimizing the displayed envelope is therefore
a precise and narrower question than finding the best possible estimator or
the actual minimum worst-case variance over the constrained graph family.

Suppose positive prices are $c_0,c_1$ and the incremental audit budget is
$C=c_0n_0+c_1n_1$. For fixed coefficients, Cauchy--Schwarz gives
\begin{equation}
 \frac14\left(\frac{a_0^2}{n_0}+\frac{a_1^2}{n_1}\right)
 \geq \frac{\left(|a_0|\sqrt{c_0}+|a_1|\sqrt{c_1}\right)^2}{4C}.
 \label{exec17:s7-constrained-budget}
\end{equation}
Equality requires $n_z$ proportional to $|a_z|/\sqrt{c_z}$. A zero coefficient
allows zero audits in that stratum; there is then no undefined sample mean in
the estimator. Continuous allocations give the general envelope solution.
Integer allocations need a feasible rounding calculation, except where the
displayed examples attain the bound exactly.

The identity $a_0-3a_1=136$ makes the remaining optimization global rather
than a search over convenient coefficients. If $c_1\leq9c_0$, then
\begin{equation}
 |a_0|\sqrt{c_0}+|a_1|\sqrt{c_1}
 \geq \frac{\sqrt{c_1}}3(|a_0|+3|a_1|)
 \geq \frac{136\sqrt{c_1}}3.
 \label{exec17:s7-stratum1-proof}
\end{equation}
Choose $\lambda=32/3$, so $a_0=0$ and $a_1=-136/3$. Auditing only stratum
one attains both inequalities. Conversely, if $c_1\geq9c_0$,
\begin{equation}
 |a_0|\sqrt{c_0}+|a_1|\sqrt{c_1}
 \geq \sqrt{c_0}(|a_0|+3|a_1|)
 \geq136\sqrt{c_0},
 \label{exec17:s7-stratum0-proof}
\end{equation}
attained by $\lambda=192$, $a_1=0$, $a_0=136$, and auditing only stratum
zero. Thus the joint minimum of the declared envelope is
\begin{equation}
 \mathcal B(C;c_0,c_1)=\frac{4624}{C}
                  \min\left(c_0,\frac{c_1}{9}\right).
 \label{exec17:s7-constrained-optimum}
\end{equation}
At $c_1=9c_0$ the two boundary designs tie. Coefficient choices between the
two kinks also attain the same bound with their corresponding allocations.
The price threshold nine is a conclusion about this envelope and information
set; it is not a universal rule for banking audits.

The stratum-one design has the particularly simple form
\begin{equation}
 \widehat v=\tfrac{58}{3}-\tfrac{136}{3}\widehat\theta_1,
 \qquad
 \widehat\theta_0=\tfrac{1-\widehat\theta_1}{3}.
 \label{exec17:s7-one-stratum}
\end{equation}
The exact pooled information supplies what would otherwise require labels
from the other stratum. This is scientifically different from assuming equal
overlap rates. A realized $\widehat\theta_1>1/2$ gives an infeasible estimated
pair even though the estimator is unbiased. Projecting this sample mean onto
$[0,1/2]$ cannot increase its distance to a true $\theta_1\in[0,1/2]$, and
therefore cannot increase the squared error of the affine value estimate.
That projected estimate need not remain unbiased. Neither operation by itself
creates a confidence interval.

At equal prices and 700 audits, assigning all 700 to stratum one gives the
envelope
\begin{equation}
 \frac{(136/3)^2}{4(700)}=\frac{1156}{1575}\simeq0.7340,
 \label{exec17:s7-pooled-count-comparison}
\end{equation}
in square dollars, compared with $28/25=1.12$ for the earlier optimized
two-mean allocation $(100,600)$. At prices $(1,4)$ and budget 1820, the
integer design $(0,455)$ gives $4624/4095\simeq1.1292$, compared with
$52/35\simeq1.4857$ for $(140,420)$. The target, sampling unit, total count
or budget, and already-available pooled information are matched. These are
exact envelope comparisons, not observed sampling improvements. Finite
enumeration on twelve small Bernoulli audit laws separately checks the
unbiased means and actual variance formula, including the zero-sample boundary
designs.~\citep{exec17_s7_extension}

\subsection{An estimated pooled constraint has a measurable error cost}
\label{exec17:s7-noisy-pool}

An actual pilot is unlikely to deliver its population pooled moment exactly.
Write $\tau=3\theta_0/4+\theta_1/4$, without fixing it at $1/4$.
Eliminating $\theta_0$ from the target gives the identity
\begin{equation}
 v_*=22-\tfrac{32}{3}\tau-\tfrac{136}{3}\theta_1.
 \label{exec17:s7-noisy-value}
\end{equation}
For unbiased estimates of both inputs, the corresponding linear value
estimate has variance
\begin{equation}
 \begin{split}
  \Var(\widehat v_*)&=\tfrac{1024}{9}\Var(\widehat\tau)
       +\tfrac{18496}{9}\Var(\widehat\theta_1)\\
       &\quad+\tfrac{8704}{9}
          \operatorname{Cov}(\widehat\tau,\widehat\theta_1).
 \end{split}
 \label{exec17:s7-noisy-variance}
\end{equation}
Separate independent audit experiments can make the covariance zero; reusing
the same stratum-one labels in the pool does not. A known pool plus one
stratum, two independent estimated summaries, and overlapping summaries are
three different observation laws. The coefficient identity does not excuse
silently replacing one law by another.

If only simultaneous bounds $|\widehat\tau-\tau|\leq d_\tau$ and
$|\widehat\theta_1-\theta_1|\leq d_1$ are warranted, the same event implies
\begin{equation}
 |\widehat v_*-v_*|\leq\tfrac{32}{3}d_\tau+
                            \tfrac{136}{3}d_1.
 \label{exec17:s7-noisy-envelope}
\end{equation}
Their joint probability must be supplied by the sampling argument; two
separately advertised nominal intervals are not automatically simultaneous.
One can intersect the joint region with the probability constraints and map
the intersection into value. Inconsistent exact equalities should instead
trigger an input review. Larger future stratum-one samples do not eliminate
fixed pool bias: ignoring a pool error of magnitude $d$ can retain a monetary
error of $32d/3$ even as those samples grow.

This comparison changes the proposed workflow in a concrete way. Before
pricing the next audit, ask which joint moments are already known, how they
were observed and what covariance or systematic error they carry. If the
pool is defensibly exact, a single strategically selected stratum can suffice
for this target. If it is noisy, preserve its error and reoptimize the joint
design. If it has an unknown truth mechanism, retain the compatible value
set and acquire different information. The source-specific missingness,
linkage and network engines can supply components only after that observation
law is declared; none can turn an uncertain aggregate into an exact scientific
constraint.~\citep{exec17_s7_extension}
